\documentclass[12pt,a4paper]{article}
\usepackage[utf8]{inputenc}
\usepackage[vietnamese]{babel}
\usepackage{amsmath}
\usepackage{amsfonts}
\usepackage{amssymb}
\usepackage{graphicx}
\usepackage{hyperref}
\usepackage{listings}
\usepackage{xcolor}
\usepackage{geometry}
\usepackage{colortbl}
\usepackage[most]{tcolorbox}

\geometry{a4paper, margin=0.8in}

\definecolor{backcolour}{rgb}{0.95,0.95,0.92}
\definecolor{outputback}{rgb}{0.1,0.1,0.1}

\newtcolorbox{expectedoutput}[1][]{
  colback=outputback,
  colupper=white,
  colframe=gray!50!black,
  title=\textbf{KẾT QUẢ DỰ KIẾN (EXPECTED OUTPUT)},
  fonttitle=\bfseries,
  enhanced,
  breakable,
  sharp corners,
  boxrule=1pt,
  #1
}

\lstdefinestyle{mystyle}{
    backgroundcolor=\color{backcolour},   
    basicstyle=\ttfamily\small,
    breaklines=true,
    frame=single,
    numbers=none
}
\lstset{style=mystyle}

\begin{document}

\begin{center}
    \Large\textbf{KỊCH BẢN DEMO BẢO VỆ ĐỒ ÁN MẠNG MÁY TÍNH NÂNG CAO} \\
    \vspace{0.3cm}
    \normalsize\textit{(Kịch bản chi tiết đạt 100\% tiêu chí Mức 2 và Mức 3)}
\end{center}

\tableofcontents
\newpage

\section{GIAI ĐOẠN 1: ĐỐI CHIẾU CẤU HÌNH VỚI TOPOLOGY}
\textit{Mục tiêu: Chứng minh thiết bị được cấu hình đúng IP và đúng lớp theo sơ đồ thiết kế.}

\subsection{1. Kiểm tra IP trên Firewall (Cửa ngõ biên)}
\textbf{Câu lệnh:} \texttt{fw ip addr show fw-eth0} và \texttt{fw ip addr show fw-eth1}
\begin{expectedoutput}
fw-eth0: 203.0.113.2/30 (Nối Internet - ext) \\
fw-eth1: 172.16.0.1/30  (Nối Backbone - core)
\end{expectedoutput}
\textbf{Lời thoại:} "Thưa thầy cô, cổng eth0 của Firewall khớp với dải IP Public nối ra ISP, cổng eth1 nối vào mạng lõi nội bộ đúng như sơ đồ."

\subsection{2. Kiểm tra IP trên các Router Lớp Distribution}
\textbf{Câu lệnh:} \texttt{dist1 ip addr show dist1-eth1}
\begin{expectedoutput}
dist1-eth1: 192.168.10.1/24 (Gateway cho Tầng 1)
\end{expectedoutput}
\textbf{Lời thoại:} "Router dist1 đóng vai trò Default Gateway cho Tầng 1 (VLAN 10). Tương tự dist2 là gateway cho Tầng 2."

\section{GIAI ĐOẠN 2: KIỂM TRA ĐỊNH TUYẾN OSPF VÀ HA (MỨC 3)}

\subsection{1. Trạng thái hội tụ OSPF}
\textbf{Câu lệnh:} \texttt{fw vtysh --vty_socket /tmp/frr_fw -c "show ip ospf neighbor"}
\begin{expectedoutput}
Neighbor ID      Pri   State           Up Time     Address         Interface
10.255.1.1         1   Full/-          05:12:44    172.16.0.2      fw-eth1
\end{expectedoutput}
\textbf{Lời thoại: } "Trạng thái \textbf{Full} chứng tỏ Firewall và Core đã trao đổi bảng định tuyến thành công."

\subsection{2. Kiểm tra bảng định tuyến tại Core}
\textbf{Câu lệnh:} \texttt{core vtysh --vty_socket /tmp/frr_core -c "show ip route ospf"}
\begin{expectedoutput}
O>* 192.168.10.0/24 [110/20] via 172.16.1.2, core-eth1 \\
O>* 192.168.20.0/24 [110/20] via 172.16.2.2, core-eth2 \\
O>* 172.16.3.0/30   [110/20] via 172.16.1.2, core-eth1 \\
\hspace*{4cm} * via 172.16.2.2, core-eth2 (ECMP)
\end{expectedoutput}
\textbf{Lời thoại: } "Core Router đã học được toàn bộ các mạng con thông qua OSPF. Đặc biệt đường HA Link (172.16.3.0) có 2 đường đi song song, chứng minh tính sẵn sàng cao."

\newpage
\section{GIAI ĐOẠN 3: KIỂM TRA BẢO MẬT VÀ NAT (MỨC 2 \& 3)}

\subsection{1. Kiểm tra NAT/PAT (Inside ra Internet)}
\textbf{Câu lệnh:} \texttt{h1 ping -c 3 203.0.113.1}
\begin{expectedoutput}
64 bytes from 203.0.113.1: icmp_seq=1 ttl=61 time=15.2 ms \\
(Ping thành công ra ISP)
\end{expectedoutput}
\textbf{Lời thoại: } "Dù h1 dùng IP Private nhưng vẫn ra được Internet nhờ cơ chế PAT (Masquerade) tại Firewall."

\subsection{2. Kiểm tra Static NAT (Từ ngoài vào DMZ)}
\textbf{Câu lệnh:} \texttt{ext curl -s 100.0.0.11:80}
\begin{expectedoutput}
<h1>Server web1 - 10.10.10.11</h1>
\end{expectedoutput}
\textbf{Lời thoại: } "External truy cập IP Public 100.0.0.11, Firewall tự động DNAT về IP nội bộ 10.10.10.11 của Web1."

\subsection{3. Kiểm chứng ACL chặn truy cập trái phép}
\textbf{Câu lệnh:} \texttt{h3 ssh 10.10.10.11}
\begin{expectedoutput}
(Lệnh bị treo - Không có phản hồi) \\
Hoặc: ssh: connect to host 10.10.10.11 port 22: Connection timed out
\end{expectedoutput}
\textbf{Lời thoại: } "Theo chính sách bảo mật, mạng Tầng 2 (h3) bị cấm SSH vào DMZ. ACL đã Drop gói tin này đúng như thiết kế."

\section{GIAI ĐOẠN 4: CÂN BẰNG TẢI VÀ TRỰC QUAN HÓA (MỨC 3)}

\subsection{1. Demo Cân bằng tải tự động}
\textbf{Câu lệnh:} \texttt{test_lb} (Trong Mininet CLI)
\begin{expectedoutput}
[SIM] Dang bom tai TCP vao PUBLIC IP 100.0.0.11... \\
Tai web1 \%: 85.2\% \\
>> CHUYEN -> web2 (Vi vuot nguong 80\%)
\end{expectedoutput}
\textbf{Lời thoại: } "Khi Web1 quá tải, script Python tự động cấu hình lại iptables để đẩy traffic sang Web2."

\subsection{2. Xem biểu đồ nhiệt Security Heatmap}
\textbf{Thao tác: } Mở ảnh \texttt{results/heatmap_acl.png}
\begin{expectedoutput}
(Hiển thị biểu đồ lưới với các mảng màu đỏ thẫm tại dải IP 192.168.20.x)
\end{expectedoutput}
\textbf{Lời thoại: } "Biểu đồ nhiệt này trích xuất từ log Firewall, cho thấy cường độ các nỗ lực xâm nhập bị chặn theo thời gian."

\end{document}
